\chapter{Integração num Segmento}\label{ch:b1:integration}

A \omterm{thm:b1:integration:def}{integral} do volume do ensino médio foi fundada em áreas tomadas
intuitivamente. Este capítulo a constrói: primeiro para \omterm{def:b1:integration:step}{funções escada},
em que a \omterm{thm:b1:integration:def}{integral} é uma soma finita, e depois para funções \omterm{def:b1:continuity:continuous}{contínuas} (e
\omterm{def:b1:continuity:continuous}{contínuas} por partes) por aproximação uniforme --- o
lugar em que o teorema de Heine (\cref{thm:b1:continuity:heine}) ganha o
seu sustento. O teorema fundamental do cálculo conecta então a
construção às primitivas, e as \omterm{thm:b1:integration:riemann}{somas de Riemann} a conectam às médias
discretas.

Ao longo do capítulo, $a < b$ são reais.

\section{Funções escada}

\begin{definition}\label{def:b1:integration:step}
$\varphi \colon \intcc{a}{b} \to \R$ é uma \emph{função
escada}\index{função escada} quando existe uma subdivisão $a = x_0
< x_1 < \dots < x_n = b$ tal que $\varphi$ é constante, igual a
$c_i$, em cada \omterm{prop:b1:reals:intervals}{intervalo} \omterm{def:b1:topology:open}{aberto} $\intoo{x_{i-1}}{x_i}$ (os valores nos
nós não são restritos). A sua \omterm{thm:b1:integration:def}{integral} é
\[
\int_a^b \varphi = \sum_{i=1}^{n} c_i\,(x_i - x_{i-1}),
\]
independente da subdivisão escolhida (refine duas subdivisões pela
sua comum: cada lado fica inalterado por refinamento).
\end{definition}

\begin{proposition}\label{prop:b1:integration:stepprops}
Nas \omterm{def:b1:integration:step}{funções escada}, a \omterm{thm:b1:integration:def}{integral} é linear, crescente ($\varphi \leq
\psi \implies \int\varphi \leq \int\psi$) e satisfaz a relação de Chasles
$\int_a^b = \int_a^c + \int_c^b$ para $a < c < b$.
\end{proposition}

\begin{proof}
O motor é a \emph{invariância por refinamento}, enunciada na
definição: inserir um nó extra $t \in \intoo{x_{i-1}}{x_i}$
numa subdivisão substitui o termo $c_i(x_i - x_{i-1})$ por
$c_i(t - x_{i-1}) + c_i(x_i - t)$ --- o mesmo número --- de modo que a
\omterm{thm:b1:integration:def}{integral} não muda por refinamento finito algum. Tome agora
$\varphi$ com subdivisão $\sigma$ e $\psi$ com subdivisão
$\sigma'$: no refinamento comum $\sigma \cup \sigma'$, as duas
são \omterm{def:b1:integration:step}{funções escada} com os \emph{mesmos} nós e, em cada peça,
$\varphi + \lambda\psi$ é constante igual a $c_i + \lambda
d_i$: a linearidade se reduz à linearidade de somas finitas. Crescimento:
$c_i \leq d_i$ em cada peça dá $\sum c_i \Delta_i \leq \sum
d_i \Delta_i$ (comprimentos $\Delta_i \geq 0$). Chasles: insira o
nó $c$ e separe a soma nele.
\end{proof}

\section{Integral de uma função contínua}

\begin{theorem}[Aproximação uniforme]\label{thm:b1:integration:approx}
Seja $f$ \omterm{def:b1:continuity:continuous}{contínua} em $\intcc{a}{b}$. Para todo $\varepsilon > 0$
existem \omterm{def:b1:integration:step}{funções escada} $\varphi, \psi$ com
\[
\varphi \leq f \leq \psi
\qquad\text{e}\qquad
\psi - \varphi \leq \varepsilon \text{ em} \intcc{a}{b}.
\]
\end{theorem}

\begin{proof}
Pelo teorema de Heine (\cref{thm:b1:continuity:heine}), $f$ é \omterm{def:b1:continuity:uniform}{uniformemente
contínua}: tome $\delta$ para $\varepsilon$ e uma subdivisão de
malha $< \delta$ (igualmente espaçada, digamos, com $n > \frac{b -
a}{\delta}$ peças). Em cada peça \omterm{def:b1:topology:closed}{fechada} $\intcc{x_{i-1}}{x_i}$, $f$
atinge um mínimo $m_i$ e um máximo $M_i$
(\cref{thm:b1:continuity:evt}), e $M_i - m_i \leq \varepsilon$
(os dois pontos extremais estão a menos de $\delta$). Defina $\varphi = m_i$
e $\psi = M_i$ em $\intoo{x_{i-1}}{x_i}$ (e $\varphi = \psi = f$
nos nós).
\end{proof}

\begin{theorem}[Definição da integral]\label{thm:b1:integration:def}
Seja $f$ \omterm{def:b1:continuity:continuous}{contínua} em $\intcc{a}{b}$. Os dois números
\[
I_-(f) = \sup\Bigl\{\int_a^b \varphi : \varphi \text{ escada},\
\varphi \leq f\Bigr\},
\qquad
I_+(f) = \inf\Bigl\{\int_a^b \psi : \psi \text{ escada},\ \psi \geq
f\Bigr\}
\]
são iguais; o seu valor comum é a
\emph{integral}\index{integral} $\int_a^b f$ (também escrita $\int_a^b
f(t)\,\dd t$). Ele coincide com a noção anterior nas \omterm{def:b1:integration:step}{funções
escada} e se estende às funções \omterm{def:b1:continuity:continuous}{contínuas} por partes, separando
$\intcc{a}{b}$ nas descontinuidades (Chasles como definição ali).
\end{theorem}

\begin{proof}
Os dois \omterm{def:b1:logic:sets}{conjuntos} são não vazios ($f$ é limitada) e toda \omterm{thm:b1:integration:def}{integral}
escada inferior é $\leq$ toda superior (monotonicidade nas \omterm{def:b1:integration:step}{funções
escada}): assim, $I_-(f) \leq I_+(f)$. Pelo \cref{thm:b1:integration:approx}, para todo
$\varepsilon$ existe um par com $\int\psi - \int\varphi \leq
\varepsilon(b - a)$: o \omterm{def:b1:reals:bounds}{supremo} e o \omterm{def:b1:reals:bounds}{ínfimo} ficam comprimidos, $I_- =
I_+$.
\end{proof}

\begin{example}[Contínua por partes, sem drama]\label{ex:b1:integration:piecewise}
A função \omterm{thm:b1:reals:floor}{parte inteira} em $\intcc{0}{3}$ é uma \omterm{def:b1:integration:step}{função escada}
disfarçada: separando nos seus saltos,
\[
\int_0^3 \lfloor t \rfloor\,\dd t
= \int_0^1 0 + \int_1^2 1 + \int_2^3 2 = 0 + 1 + 2 = 3 ,
\]
e os valores \emph{nos} pontos de salto $1, 2$ são irrelevantes:
mudar uma função em finitos pontos não muda \omterm{thm:b1:integration:def}{integral} alguma
(as \omterm{def:b1:integration:step}{funções escada} que a enquadram não são afetadas). É todo o
conteúdo da extensão a ``\omterm{def:b1:continuity:continuous}{contínua} por partes'': corte nas
finitas descontinuidades, integre cada peça \omterm{def:b1:continuity:continuous}{contínua} e
some --- Chasles como definição.
\end{example}

\begin{example}[A definição calcula, uma vez]\label{ex:b1:integration:defcompute}
Seja $f(x) = x$ em $\intcc{0}{1}$ e corte em $n$ peças iguais.
As melhores \omterm{def:b1:integration:step}{funções escada} constantes nas peças são $\varphi =
\frac{k-1}{n}$ e $\psi = \frac kn$ na $k$-ésima peça, com
\[
\int_0^1 \varphi = \sum_{k=1}^{n} \frac{k-1}{n}\cdot\frac1n
= \frac{n-1}{2n},
\qquad
\int_0^1 \psi = \sum_{k=1}^{n} \frac{k}{n}\cdot\frac1n
= \frac{n+1}{2n} .
\]
Toda \omterm{thm:b1:integration:def}{integral} inferior é $\leq I_-(f) \leq I_+(f) \leq$ toda
superior, de modo que $\frac{n-1}{2n} \leq I_-(f) \leq I_+(f) \leq
\frac{n+1}{2n}$ para todo $n$: as duas se comprimem sobre $\frac12$, e
$\int_0^1 x\,\dd x = \frac12$ sai direto da definição. A
ideia de fechamento: esta é a primeira e a última vez que integramos
pela definição --- o teorema fundamental, abaixo, substitui
todos esses cálculos por uma consulta de primitiva, o que é
todo o ponto econômico deste capítulo.
\end{example}

\begin{theorem}[Propriedades]\label{thm:b1:integration:props}
Para $f, g$ \omterm{def:b1:continuity:continuous}{contínuas} (ou \omterm{def:b1:continuity:continuous}{contínuas} por partes) em $\intcc{a}{b}$ e
$\lambda \in \R$:
\begin{enumerate}
  \item linearidade: $\int (f + \lambda g) = \int f + \lambda \int g$;
  \item monotonicidade: $f \leq g \implies \int_a^b f \leq \int_a^b g$;
        e $\bigl|\int_a^b f\bigr| \leq \int_a^b \abs f \leq (b -
        a)\, \sup\abs f$;
  \item Chasles: $\int_a^b = \int_a^c + \int_c^b$ (com a
        convenção $\int_b^a = -\int_a^b$, válida para qualquer ordem
        dos limites);
  \item positividade estrita: se $f$ é \omterm{def:b1:continuity:continuous}{contínua}, $f \geq 0$ e
        $\int_a^b f = 0$, então $f = 0$ em toda parte em
        $\intcc{a}{b}$.
\end{enumerate}
\end{theorem}

\begin{proof}
(1)--(3) passam das \omterm{def:b1:integration:step}{funções escada} ao limite pela
definição por sup/inf. A linearidade merece os detalhes uma vez: dado
$\varepsilon > 0$, enquadre $\varphi_f \leq f \leq \psi_f$ e
$\varphi_g \leq g \leq \psi_g$ com folgas $\leq \varepsilon$
(\cref{thm:b1:integration:approx}). Para $\lambda \geq 0$,
$\varphi_f + \lambda\varphi_g \leq f + \lambda g \leq \psi_f +
\lambda\psi_g$ é um enquadramento por \omterm{def:b1:integration:step}{funções escada} de folga $\leq
(1 + \lambda)\varepsilon$, e as suas integrais escada valem $\int
\varphi_f + \lambda\int\varphi_g$, etc.\
(\cref{prop:b1:integration:stepprops}): fazer $\varepsilon \to
0$ comprime $\int(f + \lambda g)$ sobre $\int f + \lambda\int g$.
Para $\lambda < 0$, multiplicar por $\lambda$ \emph{inverte} o
enquadramento de $g$ --- a \omterm{def:b1:integration:step}{função escada} inferior de $\lambda g$ é
$\lambda \psi_g$ --- e a mesma compressão roda com os papéis
trocados. A estimativa $\abs{\int f} \leq \int \abs f$ vem de
$-\abs f \leq f \leq \abs f$ e da monotonicidade.

(4) Contrapositiva: se $f(x_0) = m > 0$, a \omterm{def:b1:continuity:continuous}{continuidade} fornece um
subintervalo de comprimento $\eta > 0$ no qual $f \geq \frac m2$; a
\omterm{def:b1:integration:step}{função escada} que vale $\frac m2$ ali e $0$ no resto é $\leq f$,
de modo que $\int f \geq \frac{m\eta}{2} > 0$.
\end{proof}

\begin{example}[Chasles em ação: integrais com valores absolutos]\label{ex:b1:integration:chasleswork}
Para integrar um valor absoluto, corte onde o sinal muda.
\[
\int_0^2 \abs{x - 1}\,\dd x
= \int_0^1 (1 - x)\,\dd x + \int_1^2 (x - 1)\,\dd x
= \frac12 + \frac12 = 1 ,
\]
e, cortando $\intcc{0}{2\pi}$ em $\pi$:
\[
\int_0^{2\pi} \abs{\sin t}\,\dd t
= \int_0^{\pi} \sin t\,\dd t - \int_{\pi}^{2\pi} \sin t\,\dd t
= 2 + 2 = 4 ,
\]
ao passo que $\int_0^{2\pi} \sin t\,\dd t = 0$: o cancelamento é real,
e é por isso que o enunciado de positividade estrita
(\cref{thm:b1:integration:props}~(4)) carrega a hipótese $f
\geq 0$ --- sem ela, uma \omterm{thm:b1:integration:def}{integral} nula nada demonstra
sobre $f$. A ideia de fechamento: $\int \abs f$ mede
\emph{área}, e $\int f$ mede \emph{balanço com sinal}; a
desigualdade $\abs{\int f} \leq \int\abs f$ é o registro exato do
que o cancelamento pode destruir.
\end{example}

\section{O teorema fundamental do cálculo}

\begin{theorem}[Teorema fundamental do cálculo]\label{thm:b1:integration:ftc}
Seja $f$ \omterm{def:b1:continuity:continuous}{contínua} num \omterm{prop:b1:reals:intervals}{intervalo} $I$ e $a \in I$. A função
\[
F(x) = \int_a^x f(t)\, \dd t
\]
é de classe $C^1$ em $I$, com $F' = f$: toda função \omterm{def:b1:continuity:continuous}{contínua}
num \omterm{prop:b1:reals:intervals}{intervalo} tem primitivas. Consequentemente, para \emph{qualquer}
primitiva $G$ de $f$:\index{teorema fundamental do cálculo}
\[
\int_a^b f(t)\,\dd t = G(b) - G(a) .
\]
\end{theorem}

\begin{proof}
Fixe $x_0 \in I$ e $\varepsilon > 0$; a \omterm{def:b1:continuity:continuous}{continuidade} em $x_0$ fornece
$\delta$ com $\abs{f(t) - f(x_0)} \leq \varepsilon$ para $\abs{t -
x_0} \leq \delta$. Para $0 < \abs{h} \leq \delta$ (e $x_0 + h \in
I$), Chasles dá
\[
\frac{F(x_0 + h) - F(x_0)}{h} - f(x_0)
= \frac 1h \int_{x_0}^{x_0+h} \bigl(f(t) - f(x_0)\bigr)\dd t ,
\]
cujo valor absoluto é $\leq \frac{1}{\abs h}\cdot \abs h\,
\varepsilon = \varepsilon$ (estimativa (2), válida para qualquer ordem dos
limites). Assim, $F'(x_0) = f(x_0)$; $F' = f$ é \omterm{def:b1:continuity:continuous}{contínua}: $F$ é
$C^1$. Se $G' = f$ também, então $(G - F)' = 0$ no \omterm{prop:b1:reals:intervals}{intervalo}, de modo que $G =
F + c$ (\cref{cor:b1:derivative:monotone}), e $G(b) - G(a) = F(b)
- F(a) = \int_a^b f$.
\end{proof}

\begin{example}[Simetria antes do cálculo]\label{ex:b1:integration:parity}
Num \omterm{prop:b1:reals:intervals}{intervalo} simétrico, a paridade faz o trabalho: se $f$ é ímpar,
a substituição $t \mapsto -t$ leva $\int_{-a}^{0} f$ a
$-\int_0^a f$, de modo que
\[
\int_{-a}^{a} f(t)\,\dd t = 0 ;
\qquad\text{se $f$ é par,}\quad
\int_{-a}^{a} f = 2\int_0^a f .
\]
Assim, $\int_{-1}^{1} \frac{t^3\cos t}{1 + t^4}\,\dd t = 0$ sem primitiva à vista
(o integrando é ímpar), e
$\int_{-\pi}^{\pi} t^2\cos t\,\dd t = 2\int_0^\pi t^2\cos t\,
\dd t$. Verifique a simetria antes de recorrer a técnicas: a
\omterm{thm:b1:integration:def}{integral} mais rápida é a que nunca é calculada.
\end{example}

\begin{example}[Reconhecendo uma derivada de imediato]\label{ex:b1:integration:halfangle}
Calcule $\displaystyle\int_0^{\pi/2} \frac{\dd x}{1 + \cos x}$.
A identidade do arco metade $1 + \cos x = 2\cos^2\frac x2$ transforma o
integrando em $\frac{1}{2}\bigl(1 + \tan^2\frac
x2\bigr)$, que é exatamente a \omterm{def:b1:derivative:def}{derivada} de $\tan\frac x2$:
\[
\int_0^{\pi/2} \frac{\dd x}{1 + \cos x}
= \Bigl[\tan\frac x2\Bigr]_0^{\pi/2} = \tan\frac\pi4 = 1 .
\]
Nenhuma maquinaria de substituição foi necessária --- apenas o reflexo
de ler um integrando como a \omterm{def:b1:derivative:def}{derivada} de alguém, com o teorema
fundamental fazendo o resto. (A ferramenta sistemática por trás de tais
integrais trigonométricas, a substituição $t = \tan\frac x2$,
pertence ao instrumental padrão construído a partir do
\cref{thm:b1:integration:parts}~(2).)
\end{example}

\begin{example}[Funções definidas por integrais]\label{ex:b1:integration:definedby}
O teorema fundamental fabrica funções. Seja
\[
F(x) = \int_0^x \eu^{-t^2}\,\dd t .
\]
Nenhuma combinação de funções clássicas tem \omterm{def:b1:derivative:def}{derivada}
$\eu^{-t^2}$ (um teorema de Liouville, admitido); e, no entanto, $F$ existe,
é $C^1$ com $F'(x) = \eu^{-x^2} > 0$, estritamente crescente, ímpar
(substitua $t \mapsto -t$) e limitada: para $x \geq 1$,
\[
F(x) - F(1) = \int_1^x \eu^{-t^2}\dd t
\leq \int_1^x \eu^{-t}\dd t \leq \eu^{-1} ,
\]
de modo que $F \leq F(1) + \eu^{-1} \leq 1 + \eu^{-1}$. (O limite exato,
$\frac{\sqrt\pi}{2}$, é calculado com integrais duplas no
volume do terceiro ano.) Regra da cadeia para limites móveis: $\frac{\dd}{\dd
x}\int_x^{x^2} \eu^{-t^2}\dd t = 2x\,\eu^{-x^4} - \eu^{-x^2}$.
A ideia de fechamento: a integração \emph{cria} funções novas
a partir de antigas, com todas as suas propriedades legíveis no
integrando --- a primitiva que você não consegue escrever é ainda
uma função que você controla plenamente.
\end{example}

\begin{example}[Estimar sem avaliar]\label{ex:b1:integration:estimate}
As integrais $R_n = \int_0^1 \frac{t^n}{1 + t}\,\dd t$ não têm
forma \omterm{def:b1:topology:closed}{fechada} agradável e, no entanto, a monotonicidade as fixa com
precisão: em $\intcc{0}{1}$, $\frac12 \leq \frac{1}{1+t} \leq 1$, de modo que
\[
\frac{1}{2(n+1)} = \frac12\int_0^1 t^n\,\dd t
\;\leq\; R_n \;\leq\; \int_0^1 t^n \,\dd t = \frac{1}{n+1} :
\]
a ordem exata de decaimento ($R_n \sim$ um múltiplo de $\frac1n$,
de fato $R_n \sim \frac{1}{2n}$) com duas linhas e sem
primitiva alguma. Os problemas de fim de semana deste capítulo e do
próximo rodam exatamente sobre enquadramentos desses --- o primeiro
instinto do analista diante de uma \omterm{thm:b1:integration:def}{integral} deve ser \emph{estimá-la} e
só depois, se necessário, calculá-la.
\end{example}

\begin{example}[Valores médios]\label{ex:b1:integration:average}
A \emph{média} de uma $f$ \omterm{def:b1:continuity:continuous}{contínua} em $\intcc{a}{b}$ é
$\frac{1}{b-a}\int_a^b f$. Para o arco do seno:
\[
\frac{1}{\pi}\int_0^\pi \sin t\,\dd t
= \frac{1}{\pi}\bigl[-\cos t\bigr]_0^\pi = \frac{2}{\pi}
\approx 0.637 :
\]
um arco positivo completo tem média não $\frac12$, mas
$\frac2\pi$ --- a curva passa mais tempo alta do que um triângulo
passaria. Pelo \cref{exo:b1:integration:11} (teorema do valor médio para
integrais, $g = 1$), a média é um \emph{valor}: $\sin c =
\frac2\pi$ para algum $c \in \intoo{0}{\pi}$. E, pelas \omterm{thm:b1:integration:riemann}{somas de
Riemann} deste capítulo, a média é o limite de médias
comuns de $n$ amostras --- a ponte entre a média discreta de
dados e a média \omterm{def:b1:continuity:continuous}{contínua} de um sinal, que é como
a \omterm{thm:b1:integration:def}{integral} entra na física.
\end{example}

\begin{theorem}[Integração por partes; substituição]\label{thm:b1:integration:parts}
\begin{enumerate}
  \item Se $u, v$ são $C^1$ em $\intcc{a}{b}$:\index{integração por
        partes}
        \[
        \int_a^b u'v = \bigl[uv\bigr]_a^b - \int_a^b uv' .
        \]
  \item Se $\varphi$ é $C^1$ em $\intcc{\alpha}{\beta}$ e $f$ é
        \omterm{def:b1:continuity:continuous}{contínua} em $\varphi(\intcc{\alpha}{\beta})$:
        \index{substituição}
        \[
        \int_{\alpha}^{\beta} f\bigl(\varphi(t)\bigr)\,\varphi'(t)\,
        \dd t = \int_{\varphi(\alpha)}^{\varphi(\beta)} f(x)\, \dd x .
        \]
\end{enumerate}
\end{theorem}

\begin{proof}
(1) $(uv)' = u'v + uv'$; integre em $\intcc{a}{b}$ e aplique o
teorema fundamental à função $C^1$ $uv$.

(2) Seja $F$ uma primitiva de $f$ no \omterm{prop:b1:reals:intervals}{intervalo} imagem
(\cref{thm:b1:integration:ftc}). Então $(F \circ \varphi)' = (f
\circ \varphi)\,\varphi'$ (regra da cadeia), de modo que os dois lados valem
$F(\varphi(\beta)) - F(\varphi(\alpha))$.
\end{proof}

\begin{example}\label{ex:b1:integration:parts}
$\int_0^1 t\,\eu^t \dd t = \bigl[t\,\eu^t\bigr]_0^1 - \int_0^1
\eu^t\dd t = \eu - (\eu - 1) = 1$. E, com a substituição $x =
\sin t$ ($t \in \intcc{0}{\frac\pi2}$):
\[
\int_0^1 \sqrt{1 - x^2}\, \dd x
= \int_0^{\pi/2} \cos t \cdot \cos t \, \dd t
= \int_0^{\pi/2} \frac{1 + \cos 2t}{2}\, \dd t = \frac\pi4 ,
\]
--- um quarto do disco unitário, como a geometria exige. (A linearização
do \cref{met:b1:complex:trig} em ação.)
\end{example}

\begin{remark}[Armadilhas frequentes no cálculo integral]
(i) \emph{As substituições devem ser $C^1$ em todo o \omterm{prop:b1:reals:intervals}{intervalo}}:
a mudança $x = \frac1t$ é ilegal atravessando $0$; aplicada às cegas
a $\int_{-1}^{1}\frac{\dd x}{1 + x^2}$, ela ``demonstra'' que a
\omterm{thm:b1:integration:def}{integral} é igual ao seu próprio oposto. Quando uma substituição tem
singularidade, corte antes o \omterm{prop:b1:reals:intervals}{intervalo} (Chasles), substitua em
cada peça e só então recombine. (ii) \emph{As primitivas
logarítmicas precisam de valores absolutos}: $\int \frac{\dd x}{x - 2} =
\ln\abs{x - 2} + C$ em cada lado de $2$ separadamente --- escrever
$\ln(x - 2)$ em $\intoo{0}{1}$ é escrever o logaritmo de um
número negativo; e a constante $C$ pode diferir nos dois
lados da singularidade. (iii) \emph{Uma \omterm{thm:b1:integration:def}{integral} nula não
mata a função}: $\int_0^{2\pi}\sin = 0$; a positividade do
integrando é exigida antes de concluir que $f = 0$
(\cref{ex:b1:integration:chasleswork}). (iv) \emph{As \omterm{thm:b1:integration:riemann}{somas de Riemann}
devem ser calibradas}: em $\frac{b-a}{n}\sum f\bigl(a +
k\frac{b-a}{n}\bigr)$, o passo fora e os pontos dentro
devem corresponder à mesma subdivisão --- o erro frequente é uma soma
$\sum_{k=1}^{n} f\bigl(\frac kn\bigr)$ \emph{sem} o fator
$\frac1n$, que diverge em vez de convergir a $\int_0^1
f$. Lista de verificação antes de invocar o \cref{thm:b1:integration:riemann}:
coloque $\frac1n$ em evidência, reescreva a parcela como $f$ de $\frac kn$,
nomeie $f$ e verifique a sua \omterm{def:b1:continuity:continuous}{continuidade}.
\end{remark}

\begin{example}[Adivinhe, derive, ajuste]\label{ex:b1:integration:guessdiff}
Quanto vale $\int_1^x (\ln t)^2\,\dd t$? Adivinhe uma primitiva da
forma $t\,P(\ln t)$ com $P$ \omterm{def:b1:poly:def}{polinômio} e derive:
\[
\bigl(t\,P(\ln t)\bigr)' = P(\ln t) + P'(\ln t) .
\]
Precisamos de $P(u) + P'(u) = u^2$: tome $P(u) = u^2 - 2u + 2$
(igualando coeficientes de cima para baixo a partir de $u^2$). Portanto,
\[
\int_1^x (\ln t)^2\,\dd t
= \bigl[t\bigl((\ln t)^2 - 2\ln t + 2\bigr)\bigr]_1^x
= x(\ln x)^2 - 2x\ln x + 2x - 2 ,
\]
resultado que de outro modo se alcançaria com duas \omterm{thm:b1:integration:parts}{integrações por
partes}. A ideia de fechamento: para integrandos da forma
(\omterm{def:b1:poly:def}{polinômio} em $\ln t$) ou (\omterm{def:b1:poly:def}{polinômio} vezes $\eu^{\lambda t}$),
a primitiva tem a mesma forma --- derivar um palpite com forma
converte a integração em álgebra linear sobre coeficientes,
mais rápida e menos sujeita a erro que partes iteradas.
\end{example}

\begin{example}[A integral bumerangue]\label{ex:b1:integration:boomerang}
Calcule $I = \int_0^{\pi/2} \eu^x \cos x\,\dd x$. Integre por
partes duas vezes, derivando o fator trigonométrico a cada vez:
\[
I = \bigl[\eu^x \sin x\bigr]_0^{\pi/2}
- \int_0^{\pi/2} \eu^x \sin x\,\dd x
= \eu^{\pi/2} - J,
\qquad
J = \bigl[-\eu^x\cos x\bigr]_0^{\pi/2}
+ \int_0^{\pi/2} \eu^x\cos x\,\dd x = 1 + I .
\]
A \omterm{thm:b1:integration:def}{integral} voltou a si mesma: $I = \eu^{\pi/2} - 1 - I$,
donde
\[
I = \frac{\eu^{\pi/2} - 1}{2} .
\]
A ideia de fechamento: quando o integrando é um produto de duas
funções que se reproduzem sob derivação
($\eu^{ax}$, $\cos bx$, $\sin bx$), duas \omterm{thm:b1:integration:parts}{integrações por partes}
produzem uma equação linear para a \omterm{thm:b1:integration:def}{integral} desconhecida --- resolva-a
em vez de integrar; equivalentemente, passe por
$\eu^{(a + \iu b)x}$ (\cref{ch:b1:complex}) e tome partes
reais. Os dois caminhos dão a mesma resposta, e conferir que dão
é um teste de sanidade gratuito.
\end{example}

\section{Somas de Riemann}

\begin{theorem}[Somas de Riemann]\label{thm:b1:integration:riemann}
Seja $f$ \omterm{def:b1:continuity:continuous}{contínua} em $\intcc{a}{b}$. Então\index{somas de Riemann}
\[
S_n = \frac{b - a}{n} \sum_{k=0}^{n-1}
f\Bigl(a + k\,\frac{b-a}{n}\Bigr)
\xrightarrow[n \to \infty]{} \int_a^b f(t)\, \dd t ,
\]
e o mesmo com quaisquer pontos de avaliação dentro dos subintervalos.
\end{theorem}

\begin{proof}
$S_n$ é a \omterm{thm:b1:integration:def}{integral} da \omterm{def:b1:integration:step}{função escada} $\varphi_n$, igual a
$f(a + k\frac{b-a}{n})$ no $k$-ésimo subintervalo. Dado
$\varepsilon > 0$, a \omterm{def:b1:continuity:uniform}{continuidade uniforme} (Heine) fornece $\delta$; para
$n > \frac{b-a}{\delta}$, todo ponto de um subintervalo está a menos de
$\delta$ do seu ponto de avaliação, de modo que $\abs{f - \varphi_n} \leq
\varepsilon$ em $\intcc{a}{b}$, donde
\[
\Bigl| \int_a^b f - S_n \Bigr|
= \Bigl| \int_a^b (f - \varphi_n) \Bigr|
\leq (b-a)\,\varepsilon . \qedhere
\]
\end{proof}

\begin{omfigure}
\begin{tikzpicture}
\begin{axis}[omaxis, width=8.6cm, height=5.4cm,
    xmin=0, xmax=1.15, ymin=0, ymax=1.7,
    xtick={0,1}, ytick=\empty]
  \draw[fill=omDef!15, draw=omDef!60, thin]
    (axis cs:0,0)      rectangle (axis cs:0.125,0.4)
    (axis cs:0.125,0)  rectangle (axis cs:0.25,0.4156)
    (axis cs:0.25,0)   rectangle (axis cs:0.375,0.4625)
    (axis cs:0.375,0)  rectangle (axis cs:0.5,0.5406)
    (axis cs:0.5,0)    rectangle (axis cs:0.625,0.65)
    (axis cs:0.625,0)  rectangle (axis cs:0.75,0.7906)
    (axis cs:0.75,0)   rectangle (axis cs:0.875,0.9625)
    (axis cs:0.875,0)  rectangle (axis cs:1,1.1656);
  \addplot[omThm, very thick, domain=0:1.05, samples=60]
    {0.4 + x^2};
\end{axis}
\end{tikzpicture}

{\small Uma soma de Riemann à esquerda com $n = 8$ retângulos: quando a
malha encolhe, a \omterm{def:b1:continuity:uniform}{continuidade uniforme} força a área da escada rumo a
$\int_a^b f$.}
\end{omfigure}

\begin{example}\label{ex:b1:integration:riemannexample}
$\displaystyle\sum_{k=1}^{n} \frac{1}{n + k} = \frac 1n
\sum_{k=1}^{n} \frac{1}{1 + k/n} \xrightarrow[n\to\infty]{}
\int_0^1 \frac{\dd x}{1 + x} = \ln 2$: um limite invisível a estimativas
elementares, transparente como soma de Riemann.
\end{example}

\begin{example}[Uma segunda soma de Riemann, com calibração]\label{ex:b1:integration:riemann2}
Encontre $\lim_{n\to\infty} \sum_{k=1}^{n} \dfrac{n}{(n+k)^2}$.
Calibre:
\[
\sum_{k=1}^{n} \frac{n}{(n+k)^2}
= \frac{1}{n}\sum_{k=1}^{n} \frac{n^2}{(n+k)^2}
= \frac1n \sum_{k=1}^{n} \frac{1}{\bigl(1 + \frac kn\bigr)^2} ,
\]
uma soma de Riemann da função \omterm{def:b1:continuity:continuous}{contínua} $f(x) = \frac{1}{(1+x)^2}$ em
$\intcc{0}{1}$: o limite é
\[
\int_0^1 \frac{\dd x}{(1 + x)^2}
= \Bigl[-\frac{1}{1+x}\Bigr]_0^1 = \frac12 .
\]
A ideia de fechamento: toda a arte está na linha do meio ---
force a parcela à forma $f(\frac kn)$ ao custo de
extrair exatamente um fator $\frac1n$; uma vez certa a forma,
o teorema faz a análise e o teorema fundamental
faz a aritmética.
\end{example}

\begin{remark}[Onde a integral trabalha em seguida]
Cada construção do capítulo tem a sua sequência. As \omterm{thm:b1:integration:riemann}{somas de Riemann}
voltam no \cref{ch:b1:series} como ponte entre séries e
integrais (comparação de $\sum \frac{1}{n^\alpha}$ com $\int
\frac{\dd t}{t^\alpha}$); o resto integral é a forma mais fina
da fórmula de Taylor no \cref{ch:b1:taylor}; a definição
baseada no $\sup$ é o protótipo da \omterm{thm:b1:integration:def}{integral} de Lebesgue
do volume do terceiro ano, em que as mesmas três propriedades
(linearidade, monotonicidade, um teorema de convergência) são
reconstruídas numa classe de funções muito maior. E o problema de fim
de semana abaixo converte a \omterm{thm:b1:integration:parts}{integração por partes} em aritmética: a
irracionalidade de $\pi^2$.
\end{remark}

\section{Exercícios}

\begin{exercise}[$\star$]\label{exo:b1:integration:1}
Calcule: $\displaystyle\int_0^1 \frac{\dd x}{x^2 - 4}$
\emph{(frações parciais, \cref{ch:b1:fractions})};
$\displaystyle\int_1^{\eu} \ln t \,\dd t$;
$\displaystyle\int_0^{\pi} t \sin t\, \dd t$.
\end{exercise}

\begin{exercise}[$\star$]\label{exo:b1:integration:2}
Calcule $\displaystyle\int_0^{1} \frac{t}{(t^2+1)^2}\,\dd t$
(substituição) e
$\displaystyle\int_0^{\pi/2} \cos^3 t\, \dd t$ (escreva $\cos^3 =
\cos(1 - \sin^2)$).
\end{exercise}

\begin{exercise}[$\star$]\label{exo:b1:integration:3}
Encontre os limites, como \omterm{thm:b1:integration:riemann}{somas de Riemann}:
\[
u_n = \sum_{k=1}^{n} \frac{n}{n^2 + k^2},
\qquad
v_n = \frac{1}{n}\sqrt[n]{\frac{(2n)!}{n!\,n^n}}
\ \emph{(take logarithms)}.
\]
\end{exercise}

\begin{exercise}[$\star$]\label{exo:b1:integration:4}
Seja $f$ \omterm{def:b1:continuity:continuous}{contínua} em $\intcc{0}{1}$. Calcule $\lim_{n\to\infty}
\int_0^1 x^n f(x)\,\dd x$. \emph{(Corte $\intcc{0}{1}$ em $1 -
\delta$.)}
\end{exercise}

\begin{exercise}[$\star\star$]\label{exo:b1:integration:5}
(Cauchy--Schwarz) Para $f, g$ \omterm{def:b1:continuity:continuous}{contínuas} em $\intcc{a}{b}$, demonstre que
\[
\Bigl(\int_a^b fg\Bigr)^{\!2} \leq \int_a^b f^2 \cdot \int_a^b g^2 ,
\]
expandindo $\int_a^b (f + \lambda g)^2 \geq 0$ como um \omterm{def:b1:poly:def}{polinômio} do segundo grau em
$\lambda$. Quando há igualdade?
\end{exercise}

\begin{exercise}[$\star\star$]\label{exo:b1:integration:6}
Seja $f$ \omterm{def:b1:continuity:continuous}{contínua} em $\R$ e $T$-periódica. Demonstre que $\int_a^{a+T}
f$ não depende de $a$, e que $\frac1x \int_0^x f(t)\,\dd t
\to \frac 1T \int_0^T f$ quando $x \to +\infty$.
\end{exercise}

\begin{exercise}[$\star\star$]\label{exo:b1:integration:7}
Para $f$ \omterm{def:b1:continuity:continuous}{contínua} em $\intcc{0}{1}$ com $\int_0^1 f = \frac12$,
demonstre que $f$ tem um ponto fixo em $\intcc{0}{1}$. \emph{(Integre
$f(x) - x$ e use a positividade estrita,
\cref{thm:b1:integration:props}~(4), pela sua contrapositiva
combinada com o teorema do valor intermediário.)}
\end{exercise}

\begin{exercise}[$\star\star\star$]\label{exo:b1:integration:8}
(Integrais de Wallis) Seja $W_n = \int_0^{\pi/2} \sin^n t\,\dd t$.
\begin{enumerate}
  \item Demonstre a recorrência $n W_n = (n-1) W_{n-2}$ ($n \geq 2$) por
        partes, e calcule $W_0, W_1$, e depois $W_{2p}$ e $W_{2p+1}$
        em forma \omterm{def:b1:topology:closed}{fechada}.
  \item Demonstre que $(W_n)$ é decrescente com $\frac{W_{n+1}}{W_n}
        \to 1$; demonstre que a quantidade $(n+1)\,W_{n+1} W_n$ é
        constante, igual a $\frac\pi2$; e deduza a equivalência
        $W_n \sim \sqrt{\dfrac{\pi}{2n}}$.
\end{enumerate}
\end{exercise}

\begin{exercise}[$\star\star\star$]\label{exo:b1:integration:9}
(Niven: $\pi$ é irracional) Suponha $\pi = \frac ab$ com $a, b \in
\N^*$, e ponha, para um $n$ a ser escolhido,
\[
P(x) = \frac{x^n (a - bx)^n}{n!},
\qquad
I_n = \int_0^{\pi} P(x)\sin x\, \dd x .
\]
\begin{enumerate}
  \item Demonstre que $0 < I_n \leq \pi\,\frac{(\pi a)^n}{n!}$, que é
        $< 1$ para $n$ grande.
  \item Demonstre que $P$ e todas as suas \omterm{def:b1:derivative:def}{derivadas} assumem valores
        \emph{inteiros} em $0$ e em $\pi = \frac ab$. \emph{(Expansão
        binomial: os coeficientes de $P$ vezes $k!$ são inteiros
        para $k \geq n$; e $P(\pi - x) = P(x)$.)}
  \item Ponha $Q = P - P'' + P^{(4)} - \dots$ (uma soma finita). Verifique
        que $\bigl(Q'\sin x - Q\cos x\bigr)' = P \sin x$, e deduza
        que $I_n = Q(\pi) + Q(0)$ é um inteiro.
  \item Conclua.
\end{enumerate}
\end{exercise}

\begin{exercise}[$\star\star\star$]\label{exo:b1:integration:10}
Seja $f$ de classe $C^1$ em $\intcc{a}{b}$. Demonstre o limite do tipo
Riemann--Lebesgue
\[
\int_a^b f(t)\sin(\lambda t)\,\dd t \xrightarrow[\lambda \to
+\infty]{} 0
\]
integrando por partes. Depois demonstre-o de novo para $f$ apenas
\omterm{def:b1:continuity:continuous}{contínua}, por aproximação uniforme com \omterm{def:b1:integration:step}{funções escada}
(\cref{thm:b1:integration:approx}).
\end{exercise}

\begin{exercise}[$\star\star$]\label{exo:b1:integration:11}
(Teorema do valor médio para integrais) Sejam $f, g$ \omterm{def:b1:continuity:continuous}{contínuas} em
$\intcc{a}{b}$ com $g \geq 0$. Demonstre que existe $c \in
\intcc{a}{b}$ com
\[
\int_a^b f(t)\,g(t)\,\dd t = f(c)\int_a^b g(t)\,\dd t ,
\]
e mostre, com um exemplo, que a hipótese $g \geq 0$ não pode ser
abandonada.
\end{exercise}

\begin{exercise}[$\star\star\star$]\label{exo:b1:integration:12}
(Os momentos forçam zeros) Seja $f$ \omterm{def:b1:continuity:continuous}{contínua} em $\intcc{a}{b}$
com
\[
\int_a^b f(t)\,t^k\,\dd t = 0 \qquad \text{para} k = 0, 1,
\dots, n .
\]
Demonstre que $f$ se anula em $n + 1$ pontos distintos de
$\intoo{a}{b}$. \emph{(Se $f$ muda de sinal apenas em $z_1 < \dots
< z_m$ com $m \leq n$, integre $f$ contra $P(t) = (t - z_1)
\cdots (t - z_m)$ e use a positividade estrita.)}
\end{exercise}

\begin{remark}[Perspectivas dentro deste volume]
Três capítulos adiante se apoiam diretamente neste.
O \Cref{ch:b1:taylor} carrega o resto integral --- a
mais fina das três fórmulas de Taylor é uma \omterm{thm:b1:integration:parts}{integração por partes}
iterada $n$ vezes. O \Cref{ch:b1:series} converte o enquadramento de
somas por integrais no teste decisivo para $\sum n^{-\alpha}$,
e o seu problema de fim de semana refina esse enquadramento na
constante de Euler. O \Cref{ch:b1:curves} torna a \omterm{thm:b1:integration:def}{integral} geométrica: o
comprimento de um arco parametrizado é $\int \sqrt{x'(t)^2 +
y'(t)^2}\,\dd t$, \omterm{thm:b1:integration:def}{integral} de uma função \omterm{def:b1:continuity:continuous}{contínua} num
segmento --- precisamente o objeto aqui construído, sem necessidade de
teoria imprópria alguma. O fato mais reutilizado será o mais humilde:
$\bigl|\int f\bigr| \leq (b - a)\sup\abs{f}$, a desigualdade
que converte toda estimativa pontual numa estimativa \omterm{thm:b1:integration:def}{integral}.
\end{remark}

\section{Problema: A máquina integral de irracionalidade}

\begin{problem}[{Problema de fim de semana --- $\eu$ e $\pi^2$ são
irracionais, $\eu$ com seis casas decimais, e $\frac{22}{7} > \pi$ com
demonstração}]
\label{pb:b1:integration:1}
Um único mecanismo alimenta todo este problema: \emph{uma expressão que
deve ser um inteiro positivo e que, no entanto, é comprovadamente menor
que $1$, não pode existir}. O \Cref{exo:b1:integration:9} (Niven) o rodou uma vez para
demonstrar que $\pi \notin \Q$; aqui nós o industrializamos. A máquina
precisa de três partes: uma entrada de \emph{integralidade} (valores nas
extremidades de \omterm{def:b1:poly:def}{polinômios} bem escolhidos), uma entrada de
\emph{pequenez} (um fator $\frac{1}{n!}$ esmagando a \omterm{thm:b1:integration:def}{integral}) e uma
\emph{ponte} (a \omterm{thm:b1:integration:parts}{integração por partes}) que as conecta. Demonstramos que
$\eu$ é irracional e o calculamos com erro certificado, demonstramos o
teorema mais fino de Legendre de que $\pi^2$ é irracional, e terminamos
com a \omterm{thm:b1:integration:def}{integral} mais encantadora da análise: $\int_0^1 \frac{x^4(1 -
x)^4}{1 + x^2}\dd x = \frac{22}{7} - \pi$, que enquadra $\pi$
à mão.\index{irracionalidade!de pi ao quadrado@de
$\pi^2$}\index{irracionalidade!de e@de $\eu$}

\medskip
\noindent\textbf{Parte I --- Combustível.}
\begin{enumerate}
  \item Demonstre que $\dfrac{c^{\,n}}{n!} \to 0$ para todo $c
        > 0$ fixado \emph{(a partir de $n \geq 2c$, cada passo \omterm{def:b1:arith:divides}{divide} ao
        menos ao meio o termo)}.
  \item (Integrais beta) Demonstre, por indução em $m$ com
        \omterm{thm:b1:integration:parts}{integração por partes}:
        \[
        \int_0^1 x^{\,k}\,(1 - x)^{\,m}\,\dd x =
        \frac{k!\,m!}{(k + m + 1)!}
        \qquad (k, m \in \N).
        \]
  \item Deduza que $\displaystyle\int_0^1 \bigl(x(1-x)\bigr)^n \dd x
        = \frac{1}{(2n+1)\binom{2n}{n}}$ e --- comparando com
        a estimativa $x(1 - x) \leq \frac14$ --- a estimativa
        $\binom{2n}{n} \geq \dfrac{4^n}{2n+1}$, que corresponde a
        $\binom{2n}{n}^{1/n} \to 4$ do \cref{pb:b1:seq:1}.
  \item Demonstre o lema da pequenez usado duas vezes abaixo: para toda
        $g$ \omterm{def:b1:continuity:continuous}{contínua} e $> 0$ em $\intoo{0}{1}$,
        \[
        0 < \int_0^1 \bigl(x(1-x)\bigr)^n g(x)\,\dd x
        \leq \frac{\sup_{\intcc{0}{1}}\abs g}{4^{\,n}} .
        \]
\end{enumerate}

\medskip
\noindent\textbf{Parte II --- $\eu$: irracionalidade e depois seis
casas decimais.} Ponha $A_n = \displaystyle\int_0^1 x^n \eu^x \dd x$.
\begin{enumerate}
  \item Calcule $A_0$ e $A_1$, demonstre a recorrência $A_n = \eu
        - n A_{n-1}$ e as estimativas $0 < A_n \leq
        \dfrac{\eu}{n+1}$.
  \item Mostre por indução que $A_n = \alpha_n + \beta_n \eu$
        com $\alpha_n, \beta_n \in \Z$.
  \item Deduza que $\eu$ é irracional \emph{(se $\eu = \frac
        pq$, então $q A_n$ é um inteiro preso em
        $\intoo{0}{1}$ para $n$ grande)}. Compare com a demonstração
        do \cref{exo:b1:seq:9}: mesma conclusão, combustível diferente.
  \item Demonstre, por indução e \omterm{thm:b1:integration:parts}{integração por partes}, a fórmula
        exata do resto
        \[
        \eu = \sum_{k=0}^{n} \frac{1}{k!} + R_n,
        \qquad
        R_n = \frac{1}{n!}\int_0^1 (1 - t)^n\,\eu^{\,t}\,\dd t,
        \qquad
        \frac{1}{(n+1)!} \leq R_n \leq \frac{\eu}{(n+1)!} .
        \]
  \item Tome $n = 9$: estime $R_9$ usando $\eu < 2.75$ (a partir de
        $b_2 = 2.75$ no \cref{ex:b1:seq:e}), avalie a soma
        e conclua o enquadramento certificado $2.7182818 \leq
        \eu \leq 2.7182823$ --- seis casas decimais, $\eu \approx
        2.718282$, com demonstração.
\end{enumerate}

\medskip
\noindent\textbf{Parte III --- O teorema de Legendre: $\pi^2$ é
irracional.} Seja $f(x) = \dfrac{x^n (1 - x)^n}{n!}$, e suponha que
$\pi^2 = \frac ab$ com $a, b \in \N^*$.
\begin{enumerate}
  \item Mostre que $f(1 - x) = f(x)$ e $0 < f \leq
        \dfrac{1}{4^n\,n!}$ em $\intoo{0}{1}$.
  \item Mostre que $f^{(k)}(0)$ e $f^{(k)}(1)$ são inteiros para
        todo $k \geq 0$ \emph{(expanda $x^n(1-x)^n$ com coeficientes
        inteiros; $\frac{k!}{n!} \in \Z$ para $k \geq n$;
        depois use a simetria)}.
  \item Defina
        \[
        G = b^{\,n} \sum_{k=0}^{n} (-1)^k\,
        \pi^{2n - 2k} f^{(2k)} .
        \]
        Mostre que $G(0)$ e $G(1)$ são inteiros \emph{(cada
        $b^n \pi^{2n-2k} = a^{\,n-k}\,b^{\,k}$)}.
  \item Verifique o telescopamento $G'' + \pi^2 G = b^n \pi^{2n+2} f
        = \pi^2 a^n f$ e depois
        \[
        \frac{\dd}{\dd x}\Bigl(G'(x)\sin \pi x - \pi\,G(x)\cos
        \pi x\Bigr) = \pi^2 a^n f(x)\sin \pi x .
        \]
  \item Integre em $\intcc{0}{1}$ e conclua que
        \[
        \pi\,a^n \int_0^1 f(x)\sin(\pi x)\,\dd x = G(0) + G(1)
        \in \Z ,
        \]
        é um inteiro \emph{positivo} limitado por $\dfrac{\pi
        a^n}{4^n\,n!}$.
  \item Conclua, com a questão 1, que $\pi^2$ é irracional
        (Legendre, 1794) e que isso reforça o
        \cref{exo:b1:integration:9}: por que a irracionalidade
        de $\pi^2$ implica a de $\pi$, e não o contrário?
\end{enumerate}

\medskip
\noindent\textbf{Parte IV --- Entendendo a máquina.}
\begin{enumerate}
  \item Localize as duas forças opostas (integralidade dos dados nas
        extremidades; pequenez analítica da \omterm{thm:b1:integration:def}{integral}) e a
        ponte, nas Partes II e III. Depois explique por que o fator
        $\frac{1}{n!}$ em $f$ é o ponto crucial: se ele for removido,
        a integralidade sobrevive, mas que desigualdade morre, e para
        que frações $\frac ab$ alegadas a demonstração passa então a
        falhar?
  \item Efetividade: suponha que alguém alegue que $\pi^2 = \frac ab$
        com $a \leq 10$. Mostre que a contradição já
        aparece em $n = 7$: calcule $\pi\,(10/4)^7/7! \approx
        0.38 < 1$. A máquina não apenas
        refuta; ela refuta num estágio fixo e computável.
  \item Verifique a ponte incondicionalmente: demonstre, por duas
        \omterm{thm:b1:integration:parts}{integrações por partes}, que
        \[
        \int_0^1 x(1 - x)\sin(\pi x)\,\dd x = \frac{4}{\pi^3},
        \]
        e reconcilie com a questão 14 em $n = 1$ (mantenha
        $\pi^2$ simbólico: a identidade telescopada diz $\pi^3
        \int_0^1 f_1 \sin \pi x = -(f_1''(0) + f_1''(1)) = 4$).
  \item O que torna $\eu^x$ e $\sin \pi x$ elegíveis como
        núcleos da máquina? Identifique a propriedade (cada uma satisfaz
        uma equação diferencial linear com coeficientes
        constantes, de modo que \omterm{thm:b1:integration:parts}{integrações por partes} repetidas
        retornam ao início) e nomeie a \omterm{def:b1:topology:closure}{fronteira}: a mesma
        máquina, refinada por Hermite e Lindemann, demonstra que $\eu$
        e $\pi$ são \emph{transcendentes} --- além deste volume.
\end{enumerate}

\medskip
\noindent\textbf{Parte V --- $\frac{22}{7}$ contra $\pi$, e a
moral.}
\begin{enumerate}
  \item Estabeleça a divisão polinomial
        \[
        \frac{x^4(1-x)^4}{1 + x^2}
        = x^6 - 4x^5 + 5x^4 - 4x^2 + 4 - \frac{4}{1 + x^2},
        \]
        e deduza a célebre identidade
        \[
        \int_0^1 \frac{x^4 (1-x)^4}{1 + x^2}\,\dd x
        = \frac{22}{7} - \pi .
        \]
  \item O integrando é positivo: conclua que $\pi <
        \frac{22}{7}$. Depois, limitando $\frac{1}{1+x^2}$ entre
        $\frac12$ e $1$ e usando $\int_0^1 (x(1-x))^4 =
        \frac{1}{630}$ (questão 3), demonstre que
        \[
        \frac{22}{7} - \frac{1}{630} \;\leq\; \pi \;\leq\;
        \frac{22}{7} - \frac{1}{1260} ,
        \]
        isto é, $3.14126 \leq \pi \leq 3.14207$: duas casas
        decimais corretas, à mão.
  \item Generalize: dividindo $x^{4m}(1-x)^{4m}$ por $1 + x^2$,
        mostre que o resto é a constante $(-4)^m$ \emph{(trabalhe
        módulo $x^2 + 1$: $(1-x)^2 \equiv -2x$)}, deduza
        racionais $r_m$ com
        \[
        \abs{\pi - r_m} \leq 4^{\,1 - 5m} ,
        \]
        e verifique que $m = 1$ reproduz as questões 20--21.
  \item Confronte esses racionais com a teoria de aproximação
        do \cref{pb:b1:derivative:1}: calcule $\abs{\pi -
        \frac{22}{7}} \approx 1.26\cdot10^{-3}$ contra a
        garantia de Dirichlet $\frac{1}{49}$, e cite
        $\abs{\pi - \frac{355}{113}} \approx 2.7\cdot10^{-7}$
        contra $\frac{1}{113^2} \approx 7.8\cdot10^{-5}$:
        existem aproximações racionais excepcionalmente boas para
        $\pi$ --- coerente, pois não se sabe que $\pi$ seja
        mal aproximável.
  \item (A armadilha do inteiro, abstraída) Demonstre o lema que
        unifica tudo: \emph{se $x \in \R$ e existem
        inteiros $a_n, b_n$ com $0 < \abs{a_n + b_n x} \to 0$,
        então $x$ é irracional.} Liste as suas instâncias neste
        problema, no \cref{exo:b1:integration:9}, no
        \cref{exo:b1:seq:9} e no \cref{pb:b1:derivative:1}.
  \item Síntese, uma frase para cada: (i) as três partes da
        máquina e onde cada uma vive no instrumental deste
        capítulo; (ii) o que a \omterm{thm:b1:integration:def}{integral} contribui que o
        teorema do valor médio do \cref{pb:b1:derivative:1} não
        podia; (iii) o inventário dos resultados extraídos (duas
        irracionalidades, uma constante com seis casas decimais, um
        enquadramento de $\pi$, uma estimativa binomial); (iv) a
        \omterm{def:b1:topology:closure}{fronteira} (Hermite, Lindemann; e a mesma armadilha, rodada em
        $\zeta(2)$ e $\zeta(3)$, na aritmética do século XX).
\end{enumerate}
\end{problem}
